%% statlink-example.tex
%% A minimal worked example for the statlink package.
%%
%% Compile TWICE to resolve all tagged values:
%%   pdflatex statlink-example
%%   pdflatex statlink-example

\documentclass{article}
\usepackage{statlink}
\usepackage{booktabs}
\usepackage{hyperref}

\title{\texttt{statlink} --- Worked Example}
\author{Carlo Reggiani}
\date{}

\begin{document}

\maketitle

\noindent This document demonstrates the main features of the
\texttt{statlink} package. Values are tagged once in the tables below
and retrieved automatically in the text. Recompiling after updating
a table propagates all changes to the text without any manual editing.

%% ── Table 1: OLS regression ──────────────────────────────────────────────────

\section{OLS Regression}

Table~\ref{tab:ols} reports a simple OLS regression of log wages on
years of education and experience.

\begin{table}[h]
  \centering
  \begin{tabular}{lcc}
    \toprule
    Variable    & Coefficient & Std.\ Error \\
    \midrule
    Education   & \tagvalue{ols:education}{0.082}   & 0.011 \\
    Experience  & \tagvalue{ols:experience}{0.034}  & 0.008 \\
    Constant    & \tagvalue{ols:constant}{1.245}    & 0.103 \\
    \bottomrule
  \end{tabular}
  \caption{OLS estimates. Dependent variable: log wages.}
  \label{tab:ols}
\end{table}

An additional year of education is associated with a
\getvalue[transform=percent, round=1]{ols:education}\%
increase in wages. The return to experience is
\getvalue[transform=percent, round=1]{ols:experience}\%
per additional year.

%% ── Table 2: DID log-linear regression ───────────────────────────────────────

\section{Difference-in-Differences}

Table~\ref{tab:did} reports a difference-in-differences estimate of
the effect of a training programme on earnings, using a log-linear
specification. In this setting, the percentage effect is approximated
by $\exp(\hat{\beta}) - 1$ rather than $\hat{\beta}$ directly.

\begin{table}[h]
  \centering
  \begin{tabular}{lcc}
    \toprule
    Variable             & Coefficient & Std.\ Error \\
    \midrule
    Treatment $\times$ Post & \tagvalue{did:treatment}{0.143} & 0.032 \\
    Post                 & \tagvalue{did:post}{0.021}      & 0.018 \\
    Constant             & \tagvalue{did:constant}{6.832}  & 0.241 \\
    \bottomrule
  \end{tabular}
  \caption{DID estimates. Dependent variable: log earnings.}
  \label{tab:did}
\end{table}

The training programme increased earnings by approximately
\getvalue[transform=expminusone, round=1, percent]{did:treatment},
based on the coefficient of
\getvalue[round=3]{did:treatment}
on the interaction term.

%% ── Error handling demonstration ─────────────────────────────────────────────

\section{Error Handling}

If a key is not found, \texttt{statlink} prints a visible warning
in the text rather than failing silently:
\getvalue{nonexistent:key}.

If a non-numeric value is stored and a transform is requested,
the raw value is printed safely without crashing:
\tagvalue{note:result}{n.a.} is retrieved as
\getvalue[transform=expminusone]{note:result}.

\end{document}
